\documentclass[11pt]{article}
\usepackage[margin=1in]{geometry}
\usepackage{amsmath,amssymb,amsthm,booktabs,hyperref}
\hypersetup{colorlinks=true,linkcolor=blue,urlcolor=blue,citecolor=blue}
\newtheorem{theorem}{Theorem}
\newtheorem{lemma}[theorem]{Lemma}
\setlength{\emergencystretch}{2em}
\newcommand{\R}{\mathbb R}
\title{A sharp converse fractional-binomial inequality\\through log-convexity}
\author{Research proof and formalization packet}
\date{October 26, 2026}
\begin{document}
\maketitle
\begin{abstract}
We give a proof of the converse fractional-binomial inequality posed by
Gerhold and Simon as Conjecture~3.2. A finite sum of scaled Gamma quotients
is log-convex in the scaling parameter. Its exact values at parameters
one and two then give the desired sharp lower bound, and a stronger bound
depending on the two input weights. The Gamma-quotient convexity is
classical; we also give a finite-grid proof suited to formal verification.
A companion Lean certificate verifies the source inequality, the biased
refinement, and optimality. The prior-art review records the known
ingredients and the limitations of the novelty investigation.
\end{abstract}

\section{Statement and relation to the source question}
For $k\in\mathbb N_0$, $a>0$, and $x,y\geq0$, put
\[
 S_k(a;x,y)=\sum_{j=0}^k
 \frac{\Gamma(ak+1)}{\Gamma(aj+1)\Gamma(a(k-j)+1)}
 x^{aj}y^{a(k-j)}.
\]
We use $0^0=1$. The quotient agrees with the generalized binomial
coefficient in the primary source.

\begin{theorem}\label{thm:main}
For every $k\in\mathbb N_0$, $a\geq2$, and $x,y\geq0$,
\begin{equation}\label{eq:source}
 (x+y)^{ak}\leq 2^{a-1}S_k(a;x,y).
\end{equation}
If $x+y>0$, set $p=x/(x+y)$ and $q=y/(x+y)$. Then the stronger estimate
\begin{equation}\label{eq:biased}
 S_k(a;x,y)\geq (x+y)^{ak}
 \left(\frac{1+(q-p)^{2k}}{2}\right)^{a-1}
\end{equation}
holds. The coefficient $2^{a-1}$ in \eqref{eq:source} is optimal for each
fixed $a\geq2$, uniformly over $k,x,y$.
\end{theorem}

For $a>2$ and positive integer $k$, \eqref{eq:source} is exactly
Gerhold--Simon Conjecture~3.2 \cite{GS}, recorded as MathDB~352508.
The endpoint $a=2$, the zero cases, and the sharpness example are included
to specify the full statement; they are not separate proposed discoveries.
The biased bound is a consequence of the same argument and is likewise
part of this single proposed result.

\section{A finite-grid proof of the classical convexity ingredient}
The log-convexity of the individual scaled Gamma quotients is classical.
Alzer's complete-monotonicity theorem concerns the corresponding scaled
binomial probabilities \cite{Alzer}; Qi and Guo explicitly record their
logarithmic complete monotonicity in Section~5.2, equation~(5.1)
\cite{QG}. The following elementary proof is included to make all analytic
dependencies available to a Lean certificate.

\begin{lemma}\label{lem:grid}
If $j,l$ are positive integers and $\phi:[0,1]\to\R$ is decreasing, then
\[
 \sum_{r=1}^j\phi(r/j)+\sum_{s=1}^l\phi(s/l)
 \leq\sum_{t=1}^{j+l}\phi(t/(j+l)).
\]
Here ``decreasing'' permits equality.
\end{lemma}
\begin{proof}
Define integer ranks
\[
 L(r)=r+\lfloor lr/j\rfloor,\qquad
 R(s)=s+\lfloor(js-1)/l\rfloor.
\]
Each map is strictly increasing and takes its indicated domain into
$\{1,\ldots,j+l\}$. Their images are disjoint. Indeed, if
$L(r)=R(s)$, write $u=\lfloor lr/j\rfloor$ and
$v=\lfloor(js-1)/l\rfloor$, so $r+u=s+v$.
If $u\geq s$, then $v\geq r$ and
$js\leq ju\leq lr\leq lv<js$, a contradiction.
If $u<s$, then $v<r$ and
$lr<j(u+1)\leq js\leq l(v+1)\leq lr$, again a contradiction.
The images therefore partition $\{1,\ldots,j+l\}$ by cardinality.
Moreover,
\[
 jL(r)\leq(j+l)r,\qquad lR(s)<(j+l)s.
\]
Consequently $L(r)/(j+l)\leq r/j$ and $R(s)/(j+l)<s/l$.
Apply the decreasing property and sum over the partition.
\end{proof}

\begin{lemma}\label{lem:gamma}
For nonnegative integers $j,l$, the function
\[
 a\longmapsto
 \log\Gamma((j+l)a+1)-\log\Gamma(ja+1)-\log\Gamma(la+1)
\]
is convex on $(0,\infty)$.
\end{lemma}
\begin{proof}
If either integer is zero, the function is identically zero. Suppose both
are positive. For a positive integer $d$ and $N\geq1$ define
\begin{align*}
 A_{d,N}(a)={}&da\log(Nd)+\log((Nd)!)-Nd\log d\\
 &-\sum_{i=0}^{N-1}\sum_{r=1}^d\log(a+i+r/d).
\end{align*}
Euler's Gamma limit gives, for every $a>0$,
\[
 A_{d,N}(a)\longrightarrow\log\Gamma(da+1).
\]
This follows directly by taking logarithms in
$\Gamma(z+1)=\lim_{M\to\infty}M!M^z/\prod_{m=1}^M(z+m)$,
with $z=da$ and $M=Nd$, and grouping the denominator into blocks.
Twice differentiating the finite expression gives
\[
 A_{d,N}''(a)=\sum_{i=0}^{N-1}\sum_{r=1}^d
                   \frac{1}{(a+i+r/d)^2}.
\]
For each $i$, apply Lemma~\ref{lem:grid} to
$\phi(t)=(a+i+t)^{-2}$. It follows that
$A_{j+l,N}''-A_{j,N}''-A_{l,N}''\geq0$.
Thus $A_{j+l,N}-A_{j,N}-A_{l,N}$ is convex on $(0,\infty)$.
Passing to the pointwise limit in the defining convexity inequality proves
the assertion. No differentiation of an infinite series or integral is
needed.
\end{proof}

\section{Proof of the sharp lower bound}
First suppose $x,y>0$. Normalize $p=x/(x+y)$, $q=y/(x+y)$, so $p,q>0$ and
$p+q=1$, and define
\[
 F(a)=\sum_{j=0}^k
 \frac{\Gamma(ak+1)}{\Gamma(aj+1)\Gamma(a(k-j)+1)}
 p^{aj}q^{a(k-j)},\qquad a>0.
\]
Each summand has convex logarithm: apply Lemma~\ref{lem:gamma} with
$l=k-j$ and add the affine function $aj\log p+a(k-j)\log q$.
A finite sum of positive log-convex functions is log-convex. Explicitly,
for $0<t<1$, the termwise convexity inequalities followed by H\"older give
\[
 F((1-t)a+tb)
 \leq\sum_j f_j(a)^{1-t}f_j(b)^t
 \leq F(a)^{1-t}F(b)^t.
\]
The cases $t=0,1$ are equalities.

The ordinary binomial theorem and its even-part identity give exact anchors
\[
 F(1)=(p+q)^k=1,\qquad
 F(2)=\frac{(p+q)^{2k}+(q-p)^{2k}}2
     =\frac{1+(q-p)^{2k}}2\geq\frac12.
\]
For $a\geq2$, let $t=1/(a-1)\in(0,1]$. Since
$2=(1-t)\cdot1+t\cdot a$, convexity of $\log F$ yields
\[
 \log F(2)\leq(1-t)\log F(1)+t\log F(a)=t\log F(a).
\]
Therefore
\[
 F(a)\geq F(2)^{a-1}
       =\left(\frac{1+(q-p)^{2k}}2\right)^{a-1}
       \geq2^{1-a}.
\]
Multiplying by $(x+y)^{ak}$ proves both estimates for positive $x,y$.

For $k=0$, $S_0=1$ and both conclusions follow immediately.
For $k>0$, if one input is zero, the sole nonzero summand is the endpoint
corresponding to the other input, so $S_k=(x+y)^{ak}$; the biased factor
is then one when $x+y>0$. If both inputs vanish, both sides of
\eqref{eq:source} vanish. This covers all boundary cases.
Finally, for $k=1$ and $x=y>0$, $S_1=2x^a$ and
$(x+y)^a=2^ax^a$. Equality holds in \eqref{eq:source}, proving optimality.
\hfill$\square$

\section{Prior bounds, novelty limits, and formal verification}
The proposed contribution is the full sharp finite-sum lower bound and its
biased version, not the classical Gamma-ratio convexity. The original
source already identifies the sharp coefficient using $k=1,x=y$.
Hino--Namba's published 2026 Theorem~2.9 \cite{HN} gives, for $a\geq2$,
\[
 a^{1-a^2}\leq Z_{a,p}^{(k)}\leq a,\qquad
 Z_{a,p}^{(k)}=aF(a).
\]
The present argument gives $Z_{a,p}^{(k)}\geq a2^{1-a}$, a strictly
stronger uniform lower bound. No claim that a short deduction is novel
merely because it has been formalized is intended. The separate review
record documents the primary sources searched and any remaining overlap
risks.

The formal construction uses the same Euler-limit route as
Lemma~\ref{lem:gamma}. The rank partition, real grid comparison, finite
derivatives, Gamma limit, finite log-sum convexity, exact anchors, and
homogeneity are separate proof modules. The final source endpoint quantifies over all $k\in\mathbb N_0$, $a\geq2$, and $x,y\geq0$, with no
additional convexity or grid hypothesis. The packet's verification receipt
records all 67 declarations in seven modules, source hashes, compiler and
dependency pins, and every declaration's axiom report. The certificate uses
Lean 4.34.1 with pinned Mathlib and no proof holes or additional axioms;
the only reported axioms are \texttt{propext}, \texttt{Classical.choice},
and \texttt{Quot.sound}. Independent mathematical, source, and formal
reviews accompany the packet. No prior full sharp bound was found in the
inspected literature; this is a documented novelty screen, not an absolute
certification of worldwide priority. Alzer's complete article could not be
retrieved; its theorem and corollary types were checked through accessible
primary successors, as detailed in the review.

\begin{thebibliography}{9}
\bibitem{GS} S. Gerhold and T. Simon,
\emph{A converse to the neo-classical inequality with an application to the
Mittag-Leffler function}, Monatshefte f\"ur Mathematik 200(3) (2023), 627--645,
Conjecture~3.2. \href{https://doi.org/10.1007/s00605-022-01817-8}{DOI: 10.1007/s00605-022-01817-8}.
\bibitem{Alzer} H. Alzer,
\emph{Complete monotonicity of a function related to the binomial probability},
J. Math. Anal. Appl. 459 (2018), 10--15.
\href{https://doi.org/10.1016/j.jmaa.2017.10.077}{DOI: 10.1016/j.jmaa.2017.10.077}.
\bibitem{QG} F. Qi and B.-N. Guo,
\emph{From inequalities involving exponential functions and sums to
logarithmically complete monotonicity of ratios of gamma functions},
J. Math. Anal. Appl. 493(1), 124478 (online 2020), Section~5.2.
\href{https://pmc.ncbi.nlm.nih.gov/articles/PMC7416756/}{Open primary text}.
\bibitem{HN} M. Hino and R. Namba,
\emph{Fractional binomial distributions induced by the generalized binomial
theorem and their applications}, Math. Ann. 395, article 9 (2026), Theorem~2.9.
\href{https://doi.org/10.1007/s00208-026-03424-6}{DOI: 10.1007/s00208-026-03424-6}.
\end{thebibliography}
\end{document}
